\documentclass[11pt,a4paper]{article}

\usepackage{iftex}
\ifXeTeX
  \usepackage{fontspec}
  \setmainfont{TeX Gyre Pagella}
  \setmonofont{DejaVu Sans Mono}[Scale=0.82]
\fi

\usepackage[a4paper, top=25mm, bottom=25mm, left=28mm, right=28mm]{geometry}
\usepackage{xcolor}
\definecolor{headingblue}{RGB}{31,56,100}
\definecolor{ruleblue}{RGB}{46,90,156}
\definecolor{codebg}{RGB}{245,245,245}

\usepackage{longtable,booktabs,array,calc,multirow}
\usepackage{microtype}
\usepackage{parskip}
\setlength{\parskip}{5pt}

\usepackage{fancyvrb}
\usepackage{fvextra}
\usepackage{framed}
\newenvironment{Shaded}{%
  \setlength{\FrameSep}{4pt}%
  \definecolor{shadecolor}{RGB}{245,245,245}%
  \begin{snugshade}\footnotesize
}{\end{snugshade}}
\DefineVerbatimEnvironment{Highlighting}{Verbatim}{%
  breaklines,breakanywhere,commandchars=\\\{\}}

\newcommand{\AlertTok}[1]{\textcolor[rgb]{0.94,0.16,0.16}{#1}}
\newcommand{\AnnotationTok}[1]{\textcolor[rgb]{0.56,0.35,0.01}{\textit{#1}}}
\newcommand{\AttributeTok}[1]{\textcolor[rgb]{0.13,0.29,0.53}{#1}}
\newcommand{\BaseNTok}[1]{\textcolor[rgb]{0.00,0.00,0.81}{#1}}
\newcommand{\BuiltInTok}[1]{#1}
\newcommand{\CharTok}[1]{\textcolor[rgb]{0.31,0.60,0.02}{#1}}
\newcommand{\CommentTok}[1]{\textcolor[rgb]{0.56,0.35,0.01}{\textit{#1}}}
\newcommand{\CommentVarTok}[1]{\textcolor[rgb]{0.56,0.35,0.01}{\textbf{\textit{#1}}}}
\newcommand{\ConstantTok}[1]{\textcolor[rgb]{0.56,0.35,0.01}{#1}}
\newcommand{\ControlFlowTok}[1]{\textcolor[rgb]{0.13,0.29,0.53}{\textbf{#1}}}
\newcommand{\DataTypeTok}[1]{\textcolor[rgb]{0.13,0.29,0.53}{#1}}
\newcommand{\DecValTok}[1]{\textcolor[rgb]{0.00,0.00,0.81}{#1}}
\newcommand{\DocumentationTok}[1]{\textcolor[rgb]{0.56,0.35,0.01}{\textbf{\textit{#1}}}}
\newcommand{\ErrorTok}[1]{\textcolor[rgb]{0.64,0.00,0.00}{\textbf{#1}}}
\newcommand{\ExtensionTok}[1]{#1}
\newcommand{\FloatTok}[1]{\textcolor[rgb]{0.00,0.00,0.81}{#1}}
\newcommand{\FunctionTok}[1]{\textcolor[rgb]{0.13,0.29,0.53}{#1}}
\newcommand{\ImportTok}[1]{#1}
\newcommand{\InformationTok}[1]{\textcolor[rgb]{0.56,0.35,0.01}{\textbf{\textit{#1}}}}
\newcommand{\KeywordTok}[1]{\textcolor[rgb]{0.13,0.29,0.53}{\textbf{#1}}}
\newcommand{\NormalTok}[1]{#1}
\newcommand{\OperatorTok}[1]{\textcolor[rgb]{0.81,0.36,0.00}{\textbf{#1}}}
\newcommand{\OtherTok}[1]{\textcolor[rgb]{0.56,0.35,0.01}{#1}}
\newcommand{\PreprocessorTok}[1]{\textcolor[rgb]{0.56,0.35,0.01}{\textit{#1}}}
\newcommand{\RegionMarkerTok}[1]{#1}
\newcommand{\SpecialCharTok}[1]{\textcolor[rgb]{0.81,0.36,0.00}{#1}}
\newcommand{\SpecialStringTok}[1]{\textcolor[rgb]{0.31,0.60,0.02}{#1}}
\newcommand{\StringTok}[1]{\textcolor[rgb]{0.31,0.60,0.02}{#1}}
\newcommand{\VariableTok}[1]{\textcolor[rgb]{0.00,0.00,0.00}{#1}}
\newcommand{\VerbatimStringTok}[1]{\textcolor[rgb]{0.31,0.60,0.02}{#1}}
\newcommand{\WarningTok}[1]{\textcolor[rgb]{0.56,0.35,0.01}{\textbf{\textit{#1}}}}
\providecommand{\tightlist}{\setlength{\itemsep}{2pt}\setlength{\parskip}{0pt}}

\usepackage{titlesec}
\titleformat{\section}{\large\bfseries\color{headingblue}}{\thesection.}{0.6em}{}[\vspace{-4pt}\color{ruleblue}\rule{\linewidth}{0.6pt}]
\titleformat{\subsection}{\normalsize\bfseries\color{headingblue}}{\thesubsection}{0.6em}{}
\titleformat{\subsubsection}{\normalsize\bfseries}{\thesubsubsection}{0.6em}{}
\titlespacing{\section}{0pt}{18pt}{6pt}
\titlespacing{\subsection}{0pt}{12pt}{4pt}
\titlespacing{\subsubsection}{0pt}{10pt}{3pt}
\setlength{\tabcolsep}{6pt}
\renewcommand{\arraystretch}{1.25}

\usepackage{hyperref}
\hypersetup{
  colorlinks=true,
  linkcolor=ruleblue,
  urlcolor=ruleblue,
  pdftitle={GNU Backgammon Android Port -- MASTER v0.8.6},
  pdfauthor={clavierhaus.at},
}

\usepackage{fancyhdr}
\pagestyle{fancy}
\fancyhf{}
\fancyhead[L]{\small\color{gray}GNU Backgammon Android Port --- MASTER v0.8.6}
\fancyhead[R]{\small\color{gray}clavierhaus.at}
\fancyfoot[C]{\small\color{gray}\thepage}
\renewcommand{\headrulewidth}{0.3pt}

\title{\textbf{GNU Backgammon Android Port}}
\author{clavierhaus.at}
\date{June 2026}

\begin{document}

\maketitle\thispagestyle{fancy}

\tableofcontents\newpage

\section{GNU Backgammon Android Port --- MASTER v0.8.6}\label{gnu-backgammon-android-port-master-v085}

\textbf{GNU Backgammon by clavierhaus.at}\\
clavierhaus.at · Vienna · Austria\\
June 2026

\begin{center}\rule{0.5\linewidth}{0.5pt}\end{center}

\subsection{1. Introduction \& Project Objective}\label{introduction-project-objective}

The objective of this project is to produce a \textbf{true port} of GNU
Backgammon to the Android platform --- not a stripped-down evaluator,
but a complete port that mirrors 100\% of the functionality of the PC
version. This includes full game management, SGF import/export,
analysis, rollout, cube decisions, tutor mode, and the complete match
state machine.

The primary constraint is the monolithic nature of the original gnubg
desktop codebase, which is tightly coupled with GTK/GNOME UI components,
file I/O, global configuration state, and platform-specific threading
primitives. The approach is to compile the genuine gnubg engine source
--- every file that is portable --- and replace only what is inherently
GTK-specific with Android equivalents via a clean architectural
boundary: \texttt{android-app.c}.

\textbf{The key architectural insight:} gnubg's codebase separates
cleanly into three layers:

\begin{enumerate}
\def\labelenumi{\arabic{enumi}.}
\tightlist
\item
  \textbf{Mathematical engine} (\texttt{eval.c}, \texttt{rollout.c},
  \texttt{dice.c}, \texttt{bearoff.c} etc.) --- pure computation, zero
  UI dependencies, 100\% portable
\item
  \textbf{Game logic layer} (\texttt{play.c}, \texttt{analysis.c},
  \texttt{sgf.c}, \texttt{set.c}, \texttt{format.c} etc.) --- all GTK
  references are behind \texttt{\#if\ defined(USE\_GTK)} guards,
  portable when compiled without \texttt{USE\_GTK}
\item
  \textbf{GTK application shell} (\texttt{gnubg.c}, \texttt{progress.c},
  \texttt{gtk/*.c}) --- not portable, replaced by \texttt{android-app.c}
\end{enumerate}

Every stub that silently drops functionality is a hole in the port. The
goal is to minimise stubs to only what is genuinely GTK-specific
rendering code.

\textbf{Architecture principle:} gnubg is the authority. No game logic
is reimplemented in Kotlin or the JNI bridge. The Android layer calls
gnubg's own functions and reads \texttt{ms} (the global \texttt{matchstate})
directly. The UI reflects what the engine owns.

\subsubsection{1.1 Scope}\label{scope}

This document covers the complete build history, all architectural
decisions, every obstacle encountered and its resolution, and the exact
current state of the project as of v0.8.6. It is intended as a working
reference for any software engineer picking up or contributing to this
project.

\subsubsection{1.2 App Identity}\label{app-identity}

\begin{longtable}[]{@{}ll@{}}
\toprule\noalign{}
Property & Value \\
\midrule\noalign{}
\endhead
\bottomrule\noalign{}
\endlastfoot
App name & GNU Backgammon by clavierhaus.at \\
Android package & \texttt{com.clavierhaus.gnubg} \\
Minimum Android API & 31 (Android 12.0, Snow Cone) \\
Target ABI & \texttt{arm64-v8a} \\
Engine source & GNU Backgammon upstream (gnubg.org) \\
Repository & https://github.com/clavierhaus/gnubg-android \\
\end{longtable}

\begin{center}\rule{0.5\linewidth}{0.5pt}\end{center}

\subsection{2. Narrative of Progress}\label{narrative-of-progress}

\subsubsection{Phase 1: Monolithic Analysis \& Prototype}\label{phase-1}

Initial efforts focused on analysing the upstream GNU Backgammon source.
Attempts to build the full desktop environment proved unsustainable due
to deep dependencies on GTK, GLib, and global UI state variables
(\texttt{ms}, \texttt{ap}, \texttt{positions}). A ``Frankenstein'' build
approach established that the mathematical evaluation logic was sound in
isolation, but produced a fragile, non-portable build environment.

\subsubsection{Phase 2: Redesign --- Deterministic Engineering}\label{phase-2}

Phase 2 pivoted to a ``Clean-Room'' architecture. The upstream source
acts as a read-only reference. The \texttt{engine-core/} directory
contains only the verified, essential components. Orchestration scripts
(\texttt{pull\_dependencies.sh}, \texttt{patch\_engine.sh},
\texttt{build\_all.sh}) and the Patch Registry were established.

\subsubsection{Phase 3: The Pristine Headless Build (Completed)}\label{phase-3}

Phase 3 achieved a fully hermetic host build of the mathematical core,
confirming that the engine compiles and links cleanly on Linux when
detached from the GTK/GNOME infrastructure.

\subsubsection{Phase 4: GLib Cross-Compilation for Android (Completed)}\label{phase-4}

GLib 2.88.1 was cross-compiled for \texttt{arm64-v8a} Android API 28
from the Fedora 44 SRPM, providing the full GLib API that gnubg requires
with no behavioural deviations from upstream.

\subsubsection{Phase 5: JNI Bridge \& libgnubg-engine.so (Completed)}\label{phase-5}

The JNI bridge layer (\texttt{native-lib.c}, \texttt{stubs.c},
\texttt{Engine.kt}) and CMake build pipeline were constructed, producing
\texttt{libgnubg-engine.so} verified as ELF 64-bit ARM aarch64.

\subsubsection{Phase 6: NEON Acceleration \& Device Verification (Completed)}\label{phase-6}

ARM NEON SIMD acceleration was enabled and the engine verified on a
Pixel 8 Pro (aarch64, Android 16). Opening position evaluation and
\texttt{FindBestMove} 3-1 returning 8/5 6/5 confirmed correct.

\subsubsection{Phase 7: Cube Decisions \& Rollout (Completed)}\label{phase-7}

\texttt{Engine.cubeDecision()} and \texttt{Engine.rollout()} JNI entry
points added and verified. Multi-threaded rollout via \texttt{GThreadPool}
added in Phase 9.

\subsubsection{Phase 8: Full Game Logic Layer (Completed)}\label{phase-8}

\texttt{play.c}, \texttt{analysis.c}, \texttt{sgf.c} and all
dependencies compiled into \texttt{libgnubg-engine.so}. The
\texttt{android-app.c} architecture provides all GTK replacements.

\subsubsection{Phase 9: Multi-threaded Rollout \& SGF JNI API (Completed)}\label{phase-9}

Production-quality parallel rollout and SGF load/save JNI entry points
complete the engine feature set.

\begin{center}\rule{0.5\linewidth}{0.5pt}\end{center}

\subsection{3. Architecture}\label{architecture}

\subsubsection{3.1 Core Principle}\label{core-principle}

gnubg's \texttt{ms} (global \texttt{matchstate}) is the single source
of truth. All game state --- board, dice, turn, game status --- is owned
by gnubg. The Android layer does three things only:

\begin{enumerate}
\def\labelenumi{\arabic{enumi}.}
\tightlist
\item
  Call gnubg functions (\texttt{CommandRoll}, \texttt{CommandMove},
  \texttt{ApplySubMove}, etc.)
\item
  Read \texttt{ms} fields (\texttt{ms.anBoard}, \texttt{ms.anDice},
  \texttt{ms.fTurn}, \texttt{ms.gs}) after each call
\item
  Render what the engine reports
\end{enumerate}

No move computation, no dice tracking, no legal move determination is
performed in Kotlin. The engine exclusively determines what is legal.

\subsubsection{3.2 Key gnubg Functions Used}\label{key-gnubg-functions}

\begin{longtable}[]{@{}ll@{}}
\toprule\noalign{}
gnubg Function & Purpose \\
\midrule\noalign{}
\endhead
\bottomrule\noalign{}
\endlastfoot
\texttt{CommandNewMatch("1")} & Start a new 1-point match \\
\texttt{CommandNewGame(NULL)} & Opening roll, sets \texttt{ms.fTurn} to winner \\
\texttt{CommandRoll(NULL)} & Roll dice; for \texttt{PLAYER\_GNU} also runs \texttt{ComputerTurn} \\
\texttt{CommandMove(str)} & Apply move string, record \texttt{MOVE\_NORMAL}, call \texttt{TurnDone()} \\
\texttt{NextTurn(TRUE)} & Advance \texttt{ms.fTurn}; called inside \texttt{applyMoveString} \\
\texttt{ApplySubMove(board,iSrc,nRoll,TRUE)} & Apply one sub-move for display; no match record \\
\texttt{GenerateMoves(\&ml,board,d0,d1,fPartial)} & Generate all legal moves \\
\texttt{FormatMove(sz,board,anMove)} & Format move as gnubg point notation \\
\texttt{PipCount(board,anPip)} & Compute pip counts \\
\texttt{SwapSides(board)} & Flip \texttt{anBoard[0]} $\leftrightarrow$ \texttt{anBoard[1]} \\
\end{longtable}

\subsubsection{3.3 JNI Bridge}\label{jni-bridge}

The JNI bridge (\texttt{native-lib.c}) is a thin wrapper layer. Each
JNI function calls one or more gnubg functions and reads \texttt{ms}
directly. No logic originates in the bridge. The bridge also implements
\texttt{gnubg\_on\_board\_changed()} (called by gnubg on every board
update) with a simple LOGI call for diagnostics.

\subsubsection{3.4 \texttt{NextTurn()} and \texttt{TurnDone()}}\label{nextturn}

This is the most significant architectural subtlety. In the GTK version,
\texttt{TurnDone()} sets \texttt{fNextTurn = TRUE} and the GTK idle loop
calls \texttt{NextTurnNotify()} → \texttt{NextTurn()} to advance
\texttt{ms.fTurn}. On Android there is no GTK idle loop.

\textbf{Solution:} \texttt{applyMoveString} calls \texttt{NextTurn(TRUE)}
explicitly after \texttt{CommandMove}. \texttt{NextTurn} advances
\texttt{ms.fTurn} and triggers \texttt{ComputerTurn} if the next player
is \texttt{PLAYER\_GNU}.

\subsubsection{3.5 Board Encoding Reference}\label{board-encoding}

gnubg's \texttt{TanBoard anBoard[2][25]} as used in this port:

\begin{longtable}[]{@{}lll@{}}
\toprule\noalign{}
Array & Player & UI point mapping \\
\midrule\noalign{}
\endhead
\bottomrule\noalign{}
\endlastfoot
\texttt{anBoard[0][i]} & Engine/dark & UI point \texttt{i+1} (top, moves right to left) \\
\texttt{anBoard[1][i]} & Human/light & UI point \texttt{24-i} (bottom, moves left to right) \\
\texttt{anBoard[0][24]} & Engine checkers on bar & Bar (top half) \\
\texttt{anBoard[1][24]} & Human checkers on bar & Bar (bottom half) \\
\end{longtable}

gnubg always places the moving player's checkers in \texttt{anBoard[1]}.
When \texttt{ms.fTurn == 1} (engine moving), the board is swapped for
display via \texttt{SwapSides} to maintain the fixed human-bottom
perspective.

\begin{center}\rule{0.5\linewidth}{0.5pt}\end{center}

\subsection{4. Current JNI API}\label{jni-api}

\begin{longtable}[]{@{}ll@{}}
\toprule\noalign{}
\texttt{Engine.*} Kotlin call & JNI → gnubg mapping \\
\midrule\noalign{}
\endhead
\bottomrule\noalign{}
\endlastfoot
\texttt{initialise(weightsPath)} & \texttt{EvalInitialise}, \texttt{gnubg\_init\_tld/rollout}, \texttt{ClearMatch} \\
\texttt{newGame()} & \texttt{CommandNewMatch("1")} + \texttt{CommandNewGame(NULL)} \\
\texttt{rollDice()} & \texttt{CommandRoll(NULL)}, returns \texttt{ms.anDice} \\
\texttt{getLegalMoves(board,d0,d1,fPartial)} & \texttt{GenerateMoves} \\
\texttt{applyMoveString(str)} & \texttt{CommandMove(str)} + \texttt{NextTurn(TRUE)} \\
\texttt{applySubMove(board,iSrc,nRoll)} & \texttt{ApplySubMove(...,TRUE)} \\
\texttt{findMove(oldBoard,curBoard,d0,d1)} & \texttt{GenerateMoves} + \texttt{PositionKey} comparison + \texttt{FormatMove} \\
\texttt{formatMove(board,anMove)} & \texttt{FormatMove} \\
\texttt{getMatchBoard()} & \texttt{ms.anBoard} \\
\texttt{getMatchDice()} & \texttt{ms.anDice} \\
\texttt{getMatchTurn()} & \texttt{ms.fTurn} \\
\texttt{getMatchStatus()} & \texttt{ms.gs} \\
\texttt{pipCount(board)} & \texttt{PipCount} \\
\texttt{swapBoard(board)} & \texttt{SwapSides} \\
\texttt{evaluatePosition(board)} & \texttt{EvaluatePosition} \\
\texttt{classifyPosition(board)} & \texttt{ClassifyPosition} \\
\texttt{cubeDecision(...)} & \texttt{GeneralCubeDecisionENoLocking} + \texttt{FindCubeDecision} \\
\texttt{rollout(board,trials)} & \texttt{gnubg\_rollout} (multi-threaded) \\
\texttt{loadSGF(path)} & \texttt{CommandLoadMatch} \\
\texttt{saveSGF(path)} & \texttt{CommandSaveMatch} \\
\end{longtable}

\begin{center}\rule{0.5\linewidth}{0.5pt}\end{center}

\subsection{5. Version v0.8.6 --- Milestone}\label{version-v085}

\subsubsection{5.1 Assessment}\label{assessment-v085}

Version v0.8.6 delivers a playable one-point game with correct gnubg
move handling, sub-move interaction, and board rendering. The milestone
is less advanced than originally projected (formerly labelled v0.9.5)
due to a significant remediation effort required to remove non-gnubg game
logic that had been introduced during earlier development. The remediation
was necessary and correct: the resulting codebase is architecturally
sound.

\subsubsection{5.2 What Works}\label{what-works-v085}

\begin{itemize}
\tightlist
\item
  \textbf{Full game loop:} Opening roll → human moves → commit → engine
  rolls and moves automatically → human rolls → repeat. Both players
  complete turns correctly across multiple games.
\item
  \textbf{Opening roll:} \texttt{CommandNewGame} rolls one die per player,
  determines winner, sets \texttt{ms.fTurn} and \texttt{ms.anDice}.
  Human-wins branch enters \texttt{HUMAN\_MOVING} immediately; engine-wins
  branch enters \texttt{WAITING\_FOR\_ROLL} after engine auto-plays.
\item
  \textbf{Sub-move interaction:} Single tap on a source point applies one
  sub-move via gnubg's \texttt{ApplySubMove(board, iSrc, nRoll, TRUE)},
  mirroring the single-click behaviour of \texttt{board\_button\_release}
  in \texttt{gtkboard.c}. First die tried, then second.
\item
  \textbf{Commit:} Tap Commit button → \texttt{findMove(oldBoard,
  currentBoard, d0, d1)} locates the complete move by position key
  comparison (mirroring \texttt{update\_move()} in \texttt{gtkboard.c})
  → \texttt{CommandMove(str)} records and finalises the turn →
  \texttt{NextTurn(TRUE)} advances \texttt{ms.fTurn}.
\item
  \textbf{Undo:} Tap Undo button → display board restored to
  \texttt{oldBoard} (board at start of human turn); \texttt{remainingDice}
  reset to original roll. Pure UI state; no engine call required since
  \texttt{ApplySubMove} never modifies \texttt{ms}.
\item
  \textbf{Dice swap:} Tap dice area → swap die order in
  \texttt{remainingDice}. Mirrors clicking the dice in \texttt{gtkboard.c}.
\item
  \textbf{Board rendering:} All 24 points, bar (human and engine),
  bearoff trays, pip counts, doubling cube, dice with used-die greying.
  Human checkers always at bottom regardless of whose turn it is.
\item
  \textbf{Hit checker:} Checker hit onto bar displayed correctly in
  the bar area.
\item
  \textbf{Doubles:} Four dice displayed; all four sub-moves playable.
\item
  \textbf{Game over:} \texttt{ms.gs >= 2} detected; New Game button
  appears.
\item
  \textbf{Commit/Undo button layout:} Drawn directly on the Canvas in
  gnubg board coordinate units, positioned below the dice in the tray
  gap of the right (human) half. Gap between Commit and Undo aligns
  vertically with the gap between the two dice. The entire ensemble
  (dice top row, buttons bottom row) is centred on the mid-point of
  the right half of the board.
\end{itemize}

\subsubsection{5.3 What Does Not Yet Work}\label{what-does-not-work-v085}

\begin{itemize}
\tightlist
\item
  \textbf{Doubling cube interaction} --- cube offer, accept, pass not yet
  wired to UI
\item
  \textbf{Bearoff checker display in tray} --- checkers borne off are
  not yet rendered in the bearoff tray
\item
  \textbf{Game over screen} --- gammon/backgammon detection not yet
  surfaced in UI
\item
  \textbf{SGF save/load UI} --- JNI entry points exist but no UI
\item
  \textbf{Player panel} --- score, match length, player names not yet
  displayed
\item
  \textbf{Analysis and tutor mode} --- engine capability exists; UI
  not yet built
\item
  \textbf{16KB page alignment} --- \texttt{.so} not yet built with
  \texttt{-Wl,-z,max-page-size=16384} for Android 16 compatibility
\end{itemize}

\subsubsection{5.4 Key Lessons Learned}\label{lessons-v085}

\paragraph{gnubg source is the only authority.}

Any game logic not traceable to a specific gnubg source function is a
liability. During this development cycle, several pieces of invented
logic were introduced (cached dice, accumulated move strings, move
history stacks, manual pip-count index swaps) that each caused subtle
bugs. Every one was eventually traced to a divergence from gnubg's own
logic. The remediation rule: before writing any code, read the
corresponding gnubg source function and translate it directly.

\paragraph{The TurnDone/NextTurn split.}

\texttt{TurnDone()} does not advance \texttt{ms.fTurn}; it only sets
\texttt{fNextTurn = TRUE}. \texttt{NextTurn()} does the advancement.
In the GTK version this happens via the GLib idle loop. On Android,
\texttt{NextTurn(TRUE)} must be called explicitly after
\texttt{CommandMove}. Failure to do this leaves \texttt{ms.fTurn}
pointing at the player who just moved.

\paragraph{ms.anDice is cleared by TurnDone().}

\texttt{ms.anDice} is set by \texttt{CommandRoll} for \texttt{PLAYER\_HUMAN}
and remains valid until \texttt{TurnDone()} clears it. For
\texttt{PLAYER\_GNU}, \texttt{ComputerTurn} runs and \texttt{TurnDone()}
fires before \texttt{CommandRoll} returns. Therefore \texttt{ms.anDice}
is always 0 after an engine turn. Display dice must be tracked in
\texttt{remainingDice} (Kotlin state), not read from \texttt{ms.anDice}
after the fact.

\paragraph{ApplySubMove does not touch ms.}

\texttt{ApplySubMove(TanBoard, iSrc, nRoll, fCheckLegal)} operates
purely on the board array passed to it. It never modifies \texttt{ms}.
This makes it correct and safe for incremental display updates during
human move selection. The final \texttt{CommandMove} then commits
everything to \texttt{ms} atomically.

\paragraph{findMove via position key.}

After the human has applied sub-moves to the display board, the
complete move must be identified for \texttt{CommandMove}. The correct
method (taken directly from \texttt{update\_move()} in
\texttt{gtkboard.c}) is: \texttt{GenerateMoves} on the original board,
then compare \texttt{PositionKey(currentBoard)} against
\texttt{ml.amMoves[i].key} for each candidate. This avoids any manual
move-string construction.

\subsubsection{5.5 Immediate Next Steps (v0.8.6)}\label{next-steps-v085}

\begin{enumerate}
\def\labelenumi{\arabic{enumi}.}
\tightlist
\item
  Bearoff checker display in right tray
\item
  Gammon/backgammon detection on game over (\texttt{ms.gs} values)
\item
  Player panel: score, match length, player names
\item
  Doubling cube UI: offer, accept, pass (wired to \texttt{CommandDouble},
  \texttt{CommandTake}, \texttt{CommandDrop})
\item
  Rebuild \texttt{.so} with \texttt{-Wl,-z,max-page-size=16384} for
  Android 16 16KB page alignment
\end{enumerate}

\begin{center}\rule{0.5\linewidth}{0.5pt}\end{center}


\subsection{6. Version v0.8.6 --- Engine Dice, Roll Button, Gameplay Stability}\label{version-v086}

\subsubsection{6.1 Milestone}\label{milestone-v086}

Version v0.8.6 completes the core visual game loop. Engine dice are now
persistently displayed after each engine move. The Roll button appears on
the board canvas in the right half, aligned with where the player's dice
will appear. All dice and buttons are drawn in a single coordinate system
(board units) with tap detection in the same system, guaranteeing
pixel-perfect alignment on any device.

\subsubsection{6.2 Delivered}\label{delivered-v086}

\paragraph{Engine Dice Display}

Engine dice were previously invisible after each move because
\texttt{ms.anDice} is cleared by \texttt{TurnDone()} immediately after
the engine plays. The correct source is the move record:
\texttt{getMoveRecordDice()} reads \texttt{plLastMove->p->anDice[0/1]},
which persists until the next move record is written. This mirrors
exactly what gnubg's GTK board does for display.

Engine dice are rendered in the left half of the board, centred on the
same horizontal line as the cube and the player's Roll button. Color:
\texttt{Color.White.copy(alpha = 0.3f)} — same as spent player dice.

\paragraph{Roll Button on Canvas}

The Roll button is drawn directly on the Canvas in board-unit coordinates,
centred in the right half of the board. Tap detection uses the same
board-unit coordinate system as all other board elements. This replaces
the left-panel Roll button and makes the interaction spatially natural.

\paragraph{fAutoRoll = FALSE}

\texttt{fAutoRoll} was defined as \texttt{TRUE} in \texttt{android-app.c}
line 494. With \texttt{fAutoRoll = TRUE}, \texttt{NextTurn()} automatically
calls \texttt{CommandRoll(NULL)} for the human player immediately after the
engine moves, bypassing the Roll button entirely. Setting
\texttt{fAutoRoll = FALSE} restores correct behaviour.

\paragraph{Commit Re-entry Guard}

Multiple rapid taps on the Commit button queued multiple coroutines, each
calling \texttt{applyMoveString} on the same stale state. Fixed by
setting \texttt{phase = ENGINE\_THINKING} immediately after
\texttt{applyMoveString} succeeds, blocking any subsequent \texttt{confirm()}
or Roll tap until the engine completes its turn and \texttt{readMatchState}
transitions to \texttt{HUMAN\_MOVING} or \texttt{WAITING\_FOR\_ROLL}.

\paragraph{Used Dice Graying}

Player dice that have been consumed by a sub-move are rendered at
\texttt{Color.White.copy(alpha = 0.3f)}. Both dice are always shown
(using \texttt{originalDice} for display); \texttt{usedCount =
diceToShow.size - remainingDice.size} drives the graying index.

\paragraph{Bearoff Tray Checkers}

Borne-off checkers are computed as \texttt{15 - sum(board[25..49])}
(human) and \texttt{15 - sum(board[0..24])} (engine), then rendered
as stacked checker slots in the respective tray areas.

\paragraph{Winner Detection}

Previously inverted. \texttt{getMatchWinner()} reads
\texttt{ms.anScore[0]} vs \texttt{ms.anScore[1]}, updated by
\texttt{ApplyGameOver()} in \texttt{play.c}. Returns 0 for human win,
1 for engine win.

\paragraph{Bar Re-entry}

Tap on the bar area (bottom half of the centre bar) sets \texttt{src=24}
for \texttt{ApplySubMove}. Previously the bar had no tap handler and
checkers on the bar could not be re-entered.

\paragraph{No Legal Moves}

When \texttt{CommandRoll} finds no legal moves, it calls
\texttt{TurnDone()} internally. \texttt{PlayMove} with an empty move
sets \texttt{ms.fTurn = !fPlayer}. The fix: \texttt{while(fNextTurn)
NextTurn(TRUE)} in the \texttt{rollDice} JNI function, mirroring the
non-GTK loop in \texttt{play.c} line 3596.

\subsubsection{6.3 Architecture Notes}\label{arch-v086}

The single most important lesson from this phase: \textbf{every piece of
display state must be read from gnubg's own data structures, not
reconstructed in Kotlin}. The engine dice problem was a direct consequence
of reading \texttt{ms.anDice} (transient) instead of the move record
(persistent). The move record is the authoritative source for everything
that has happened in the game.

\subsubsection{6.4 Immediate Next Steps (v0.8.7)}\label{next-steps-v086}

\begin{enumerate}
\def\labelenumi{\arabic{enumi}.}
\tightlist
\item
  Max 5 checkers visual — stack overflow display with count in top checker
\item
  Doubling cube UI — offer, accept, pass (\texttt{CommandDouble},
  \texttt{CommandTake}, \texttt{CommandDrop})
\item
  Player panel — score, match length, player names
\item
  Gammon/backgammon detection on game over
\item
  Rebuild \texttt{.so} with \texttt{-Wl,-z,max-page-size=16384} for
  Android 16 16KB page alignment
\end{enumerate}

\begin{center}\rule{0.5\linewidth}{0.5pt}\end{center}
\emph{GNU Backgammon Android Port --- MASTER v0.8.6 --- clavierhaus.at
--- June 2026}

\end{document}
